\documentclass[12pt,a4paper,UTF8]{ctexart}
\usepackage{geometry}
\geometry{left=2.5cm,right=2.5cm,top=2.5cm,bottom=2.5cm}
\usepackage{amsmath,amssymb}
\usepackage{array}
\usepackage{multirow}
\usepackage{makecell}
\usepackage{booktabs}
\usepackage{color,xcolor}
\usepackage{colortbl} % 修复表格颜色报错
\usepackage{mhchem}   % 修复化学式报错

\begin{document}

\begin{center}
    \Huge\bfseries\textcolor{red!80!black}{期末报告的说明}
\end{center}

\vspace{0.5cm}
\noindent\textbf{主题：}与新能源应用和开发相关的前沿话题\\
\textbf{时长：}15~20分钟（包括问答环节）\\
\textbf{时间：}第9周$\sim$第10周（4月27日、4月29日、5月6日）\\
\textbf{顺序：}待定（以4人小组的形式进行）

\vspace{0.5cm}

\noindent\textbf{\textcolor{purple!80!black}{前沿技术与关键材料类}}
\begin{enumerate}
    \item 钙钛矿太阳电池
    \item 钠离子电池
    \item 全固态电池
    \item 超临界\ce{CO2}发电
    \item 光裂解水制氢
    \item 绿色甲醇与绿色氨
    \item 可持续航空燃料
    \item 超热岩地热能
    \item 小型模块化反应堆
    \item \ce{CO2}捕集、利用与封存
\end{enumerate}

\vspace{0.3cm}
\noindent\textbf{\textcolor{purple!80!black}{政策治理与国际竞争}}
\begin{enumerate}
    \item 碳交易与碳税
    \item 绿电交易与绿证
    \item 欧盟CBAM机制
    \item 关键矿产供应链
    \item 关税、补贴和产业政策
    \item 绿氢认证和碳足迹
    \item 核电复兴的政策逻辑
    \item 能源转型中，如何面对煤电地区与传统产业
    \item 大型能源项目审批为何越来越受到生态约束
\end{enumerate}

\vspace{0.3cm}
\noindent\textbf{\textcolor{purple!80!black}{市场产业与商业模式类}}
\begin{enumerate}
    \item 光伏的产能过剩与价格战
    \item 储能项目靠什么赚钱
    \item 虚拟电厂的商业模式
    \item 家庭光储充一体化
    \item 电池回收与梯次利用
    \item 绿氢项目为什么停在示范和规划阶段
    \item 数据中心通过绿电采购和PPA会重塑哪些市场关系
    \item 新能源就业与技能缺口
    \item 地缘政治与融资环境
\end{enumerate}

\vspace{0.5cm}
\noindent\textbf{内容要求：}
\begin{enumerate}
    \item \textbf{封面：}题目、课程名称、小组成员、学号
    \item \textbf{选题与核心问题：}具体在调研什么问题
    \item \textbf{背景与现实意义：}说明这个问题为什么重要，为什么是当前的前沿，而不是一般性的旧话题
    \item \textbf{基本概念与原理机制：}
          \begin{itemize}
              \item 技术类题目讲工作原理或技术路线
              \item 政策类题目讲制度设计、目标和作用机制
              \item 市场类题目讲产业链、成本结构或盈利逻辑
          \end{itemize}
    \item \textbf{最新进展与关键数据：}必须有数据支撑，至少放两张带来源的图表，讲清楚发展现状，而不是只写大段文字
    \item \textbf{典型案例：}至少分析一个真实案例，还可以进行比较分析
    \item \textbf{主要挑战或争议}
    \item \textbf{未来趋势与小组判断}
    \item \textbf{结论}
\end{enumerate}

\vspace{0.5cm}
\noindent\textbf{评分标准：}
\begin{center}
\resizebox{\textwidth}{!}{%
\begin{tabular}{|c|c|c|}
\hline
\rowcolor{blue!15} \textbf{维度} & \textbf{分值} & \textbf{重点看什么} \\
\hline
选题与问题意识 & 10 & 题目是否具体，是否切合课程，是否具有当前前沿性 \\
\hline
资料质量与事实准确性 & 15 & 来源是否可靠，数据是否有出处，是否存在明显事实错误 \\
\hline
分析深度与比较能力 & 20 & 是否有比较、归纳、判断 \\
\hline
结构逻辑与时间控制 & 15 & 开头是否清楚，中间是否连贯，结论是否回扣问题 \\
\hline
PPT与图表表达 & 10 & 版面是否清晰，图表是否易懂，文字是否过多 \\
\hline
总结与未来判断 & 10 & 结论是否明确，是否提出有依据的趋势判断 \\
\hline
口头表达 & 10 & 发音、语速、清晰度、自然度 \\
\hline
问答与团队协作 & 10 & 回答问题是否准确，组内分工是否均衡，每位成员是否有参与 \\
\hline
\end{tabular}%
}
\end{center}

\end{document}